% Slide 2 – problem and target group (criterion: relevance and usefulness for the group, 25 %).
\begin{frame}{“Accessible / not accessible” is not enough}
  \begin{columns}[T,onlytextwidth]
    \begin{column}{0.49\textwidth}
      \begin{block}{Problem}
        \begin{itemize}
          \item One label does not say whether I can get in with a wheelchair, whether the door is 90~cm wide,
                or whether there is a toilet.
          \item Information online has no source and no date: a broken lift turns “accessible” into a wasted trip.
          \item Missing information is shown as if it were confirmation, and the City has no capacity to
                maintain a database by hand.
        \end{itemize}
      \end{block}
    \end{column}
    \begin{column}{0.49\textwidth}
      \begin{block}{Who it is for -- scope of the prototype}
        \begin{itemize}
          \item \textbf{Primary group:} wheelchair users and parents with prams.
          \item The same data serves further profiles: Low vision, Senior, After treatment, Luggage, Dog, Bicycle.
          \item Residents and tourists (in Polish, English, German and Ukrainian) and \textbf{venue owners},
                who get their own panel.
        \end{itemize}
      \end{block}
    \end{column}
  \end{columns}
  \vspace{1mm}
  \begin{exampleblock}{Our answer}
    Accessly answers \textbf{“will this place work for me?”} with concrete facts, each with a
    \textbf{source, a date and a verification status}. No data means no data, never accessibility.
  \end{exampleblock}

  \note[item]{Open with the story: Anna, a wheelchair user, wants a coffee in the centre. Last year's map says “accessible”; today there is a step at the door and nobody knows whether the toilet is downstairs.}
  \note[item]{Primary group from the brief: wheelchairs and prams (the “Wheelchair” and “Pram” profiles). The other profiles reuse the same attributes, so they cost nothing extra.}
  \note[item]{Stress the two groups: end users and owners; without owners the data will not stay current.}
  \note[item]{A profile is a set of preferences (steps, lift, quiet), not disability information; this comes back on the privacy slide.}
\end{frame}
